% Capítulo 2: Estado del Arte

\section{El problema del olvido catastrófico}

El olvido catastrófico fue caracterizado por primera vez de forma rigurosa por McCloskey y
Cohen~\cite{ref1} y Ratcliff~\cite{ref7}, que demostraron experimentalmente cómo incluso un
entrenamiento mínimo sobre nueva información borraba casi por completo las representaciones
previamente aprendidas en modelos conexionistas. A diferencia del olvido humano, que es gradual y
selectivo, el olvido en redes artificiales es abrupto y masivo: basta con actualizar los pesos
mediante descenso de gradiente sobre un nuevo conjunto de datos para que el rendimiento en los datos
anteriores caiga drásticamente.

La raíz del problema reside en la naturaleza distribuida de las representaciones en redes
neuronales: los mismos pesos participan en la codificación de múltiples patrones, de modo que su
modificación para adaptarse a los nuevos datos necesariamente interfiere con los antiguos. Este
conflicto entre capacidad de aprender cosas nuevas y capacidad de retener las ya aprendidas se
denomina el \textbf{dilema estabilidad-plasticidad}~\cite{ref8}.

\section{Taxonomía de soluciones}

Las estrategias para mitigar el olvido catastrófico se agrupan en tres grandes familias~\cite{ref9}:

\subsection{Métodos basados en regularización}

Estos métodos añaden un término de penalización a la función de pérdida para evitar que los pesos
más importantes para las tareas anteriores cambien demasiado. El más influyente es \textbf{EWC}
(\textit{Elastic Weight Consolidation}, Kirkpatrick et al., 2017)~\cite{ref4}, que usa la diagonal
de la matriz de información de Fisher como medida de la importancia de cada peso. Formalmente, la
función de pérdida para la tarea B es:

\begin{lstlisting}
L_B(theta) = L_B_orig(theta) + (lambda/2) * Sum_i F_i * (theta_i - theta*_A,i)^2
\end{lstlisting}

donde $\theta^*_A$ son los pesos óptimos para la tarea A, $F_i$ es la importancia del peso $i$, y
$\lambda$ controla el equilibrio entre plasticidad y estabilidad. \textbf{SI} (\textit{Synaptic
Intelligence}, Zenke et al., 2017)~\cite{ref5} estima las importancias de forma online durante el
entrenamiento, evitando el coste computacional del cálculo de la Fisher.

El principal inconveniente de estos métodos es que la penalización es solo una aproximación: no
garantizan la ausencia total de olvido, y la importancia calculada sobre una tarea puede no ser
compatible con las importancias calculadas sobre múltiples tareas sucesivas.

\subsection{Métodos basados en repetición (\textit{replay})}

Estos métodos mantienen una memoria de ejemplos pasados (reales o generados) que se mezclan con los
datos nuevos durante el re-entrenamiento. El \textbf{replay exacto} almacena directamente
subconjuntos de datos históricos (\textit{coreset methods}), con el coste de violar potencialmente
restricciones de privacidad o almacenamiento. El \textbf{replay generativo} (Shin et
al., 2017)~\cite{ref10} usa un modelo generativo (típicamente una GAN) entrenado sobre las tareas
pasadas para sintetizar ejemplos artificiales que se mezclan con los datos nuevos.

El replay generativo elimina la necesidad de almacenar datos reales, pero introduce un nuevo modelo
que debe entrenarse a su vez de forma continua, trasladando el problema del olvido al generador.

\subsection{Métodos basados en arquitectura}

Estos métodos modifican la estructura de la red para aislar el conocimiento de cada tarea. Las
\textbf{Redes Neuronales Progresivas} (PNN, Rusu et al., 2016)~\cite{ref6} añaden una columna nueva
a la red para cada tarea nueva, con conexiones laterales a las columnas anteriores, pero congelando
los pesos previos. Esto garantiza que las tareas antiguas no se olvidan, pero el tamaño del modelo
crece linealmente con el número de tareas.

Los métodos basados en máscaras binarias (\textit{PackNet}, Mallya y Lazebnik, 2018~\cite{ref11};
\textit{HAT}, Serrà et al., 2018~\cite{ref12}) reservan subconjuntos de pesos para cada tarea, pero
requieren conocer el número de tareas de antemano y tienen una capacidad finita.

\section{Linealidad local de redes ReLU}

Una propiedad matemática fundamental de las redes con activación ReLU es que definen funciones
\textbf{continuas y lineales a tramos}: para una entrada dada, la composición de transformaciones
afines y funciones ReLU resulta en una transformación lineal afín en la región convexa del espacio
de entrada que contiene a esa muestra~\cite{ref3}. Esta región queda determinada por el
\textbf{patrón de activación} binario: el vector que indica qué neuronas tienen salida positiva y
cuáles cero.

Esta propiedad ha sido estudiada desde la perspectiva de la capacidad representacional (las redes
ReLU profundas pueden generar un número de regiones lineales exponencial con la
profundidad~\cite{ref3}) y más recientemente desde la perspectiva del análisis de
robustez~\cite{ref13}. Qsimov explota esta propiedad de forma constructiva: en lugar de analizar
las regiones lineales de la red completa, extrae y enumera los \textbf{caminos de conexión activos}
entre capas, que son la representación dual de las regiones lineales cuando se consideran las
conexiones capa a capa.

\section{Sistemas lineales para re-entrenamiento}

La idea de reformular el entrenamiento de redes neuronales como un problema de mínimos cuadrados es
clásica en el contexto de las capas de salida. La técnica \textbf{LWTA} (\textit{Local Winner Takes
All}, Srivastava et al., 2013~\cite{ref14}) y los métodos de \textbf{extreme learning machines}
(ELM, Huang et al., 2006~\cite{ref15}) fijan los pesos de las capas ocultas aleatoriamente y
entrenan solo la última capa con mínimos cuadrados. Qsimov se distingue de estos enfoques en que
los pesos del modelo base no son aleatorios sino pre-entrenados, y en que la representación lineal
no se aplica solo a la última capa sino a una secuencia de capas cuyo punto de inicio es
configurable.

\section{Resumen y posicionamiento de Qsimov}

La tabla siguiente sitúa Qsimov en relación con las principales aproximaciones al aprendizaje
continuo:

\begin{table}[htbp]
  \centering
  \caption{Comparativa de métodos de aprendizaje continuo.}
  \label{tab:comparativa}
  \resizebox{\textwidth}{!}{%
  \begin{tabular}{lcccc}
    \toprule
    \textbf{Método} & \textbf{Sin datos previos} & \textbf{Sin arq. creciente} & \textbf{Garantía no-olvido} & \textbf{Velocidad} \\
    \midrule
    Fine-tuning estándar            & \checkmark & \checkmark & $\times$              & Alta \\
    EWC~\cite{ref4}                 & \checkmark & \checkmark & Parcial               & Media \\
    Replay exacto                   & $\times$   & \checkmark & Parcial               & Baja \\
    Replay generativo~\cite{ref10}  & \checkmark & $\times$   & Parcial               & Baja \\
    PNN~\cite{ref6}                 & \checkmark & $\times$   & \checkmark            & Media \\
    \textbf{Qsimov (sist. lineal)}  & \checkmark & \checkmark & \checkmark (datos vis.) & Muy alta \\
    \bottomrule
  \end{tabular}}%
\end{table}

Qsimov es el único método en esta comparativa que combina ausencia de necesidad de datos históricos,
tamaño de modelo constante, garantía exacta de no-olvido para los datos incluidos, y actualización
de coste prácticamente constante independientemente del histórico acumulado.

El hueco que ocupa Qsimov en el panorama actual es, por tanto, la intersección de tres propiedades
que hasta ahora no se daban simultáneamente en ningún método revisado: no necesitar acceso a datos
pasados, no aumentar la arquitectura y ofrecer garantía formal de no-olvido. El presente trabajo
evalúa empíricamente si estas propiedades teóricas se traducen en ventajas medibles en escenarios de
referencia estándar.
